%% ma: L10-B13-0113
%% lop: 10
%% chuong: 5
%% bai: 13
%% dang: 10.13.12
%% loai: TLN
%% muc: NB
%% dap_an: "7"
%% nguon: "chuyen-de-cac-so-dac-trung-cua-mau-so-lieu-khong-ghep-nhom-toan-10-chuong-trinh-moi.pdf (GV. Phan Nhật Linh) (P03), Bài 02 Dạng 2 Phần III Câu 4"
%% kien_thuc: "Lớp 10 Bài 13 (mốt)"
%% trang_thai: chinh-thuc
%% ghi_chu: "Cùng số liệu Phần I Câu 11 nhưng khác loại câu. Tự duyệt theo yêu cầu giáo viên 22/10/2026."
\begin{ex}
Để khảo sát kết quả môn Toán (thang điểm 10) trong kì thi tuyển sinh năm vừa qua của trường A, người ta chọn một mẫu gồm $100$ học sinh. Điểm môn Toán của các học sinh được cho ở bảng tần số sau:
\begin{center}
\begin{tabular}{|c|c|c|c|c|c|c|c|c|c|c|c|}\hline
Điểm & $0$ & $1$ & $2$ & $3$ & $4$ & $5$ & $6$ & $7$ & $8$ & $9$ & $10$\\\hline
Tần số & $1$ & $1$ & $3$ & $5$ & $8$ & $13$ & $19$ & $24$ & $14$ & $10$ & $2$\\\hline
\end{tabular}
\end{center}
Tìm mốt của mẫu số liệu.
\shortans{7}
\loigiai{Giá trị $7$ có tần số lớn nhất ($24$) nên $M_o=7$.}
\end{ex}
